\section{ITS Identity Theft Data Overview}
\label{app:demographics}


\subsection{Incident Characteristics}
\label{sec:incident_characteristics}

Next, summaries of the incidents affecting these victims were examined. The averages in Table \ref{tab:theft_summary_avg} represent the mean of the weighted survey values across the six survey years. It is important to note that for the "Type of Identity Theft" analysis, the percentages are calculated as a share of the total victim population for a given year. Because a single individual can experience multiple forms of identity theft (e.g., both credit card misuse and new account fraud), these categories are not mutually exclusive, and their percentages are not expected to sum to 100\%. Similarly, the percentages for the "How Misuse Was Discovered" and "How Information Was Obtained" sections do not total 100\%. This is because many respondents did not know the answer, and due to the skip logic of the survey. The questionnaire uses skip patterns, where a respondent's answer to one question determines whether they are asked subsequent, more detailed questions. For example, the question detailing the specific method of theft was only asked if a victim indicated that they knew how their information was obtained. If they answered "no," they were skipped past the follow up, and thus are not represented in any of those categories.

\begin{table}[h!]
  \centering
  \caption{Average Summary of Identity Theft Incidents (2008-2021).}
  \label{tab:theft_summary_avg}
  \begin{tabular}{lrr}
    \toprule
    \textbf{Category} & \textbf{N} & \textbf{\% of victims} \\
    \midrule
    \multicolumn{3}{l}{\textbf{Type of Identity Theft\tnote{*}}} \\
    \quad Existing Bank Account Misuse & 8,290,072 & 41.77 \\
    \quad Existing Credit Card Misuse & 9,476,109 & 48.23 \\
    \quad Other Existing Account Misuse & 2,598,402 & 12.28 \\
    \quad New Account Fraud & 1,443,137 & 7.82 \\
    \quad Other Misuse of Personal Information & 1,088,330 & 5.39 \\
    \quad Multiple Types & 2,734,282 & 13.76 \\
    \midrule
    \multicolumn{3}{l}{\textbf{How Misuse Was Discovered\tnote{**}}} \\
    \quad Noticed Problem with Account & 6,698,163 & 33.65 \\
    \quad Notified by Financial Institution & 7,790,053 & 37.69 \\
    \quad Notified by Other Person/Organization & 1,852,203 & 8.76 \\
    \quad Received Bill/Item Not Ordered & 679,982 & 3.63 \\
    \quad Checked Credit Report & 140,646 & 0.81 \\
    \quad Other/Unspecified & 1,807,364 & 8.37 \\
    \quad Unanswered/Not Applicable & 1,518,891\tnote{***} & 7.09 \\
    \midrule
    \multicolumn{3}{l}{\textbf{How Information Was Obtained\tnote{**}}} \\
    \quad Lost or Stolen Physical Item & 932,295 & 5.10 \\
    \quad Stolen During Transaction & 2,403,504 & 12.12 \\
    \quad Hacking/Computer Theft & 278,239 & 1.38 \\
    \quad Deceived by Scam/Phishing & 227,154 & 1.12 \\
    \quad Data Breach (Company/Employer) & 746,451 & 3.73 \\
    \quad Other & 574,991 & 2.67 \\
    \quad Unanswered/Not Applicable & 16,108,164\tnote{***} & 73.88 \\
    \bottomrule
  \end{tabular}
  % \begin{tablenotes}
  %   % \item[*] Categories are non-exclusive and may sum to more than 100\% as a single victim can experience multiple types of identity theft.
  %   % \item[**] The sum of percentages in these sections is approximately 100\% of the victim population.
  %   % \item[***] The Unanswered/Not Applicable row aggregates the weighted counts for respondents who were not asked the question due to survey skip logic or who had a missing data code (e.g., 'Don't Know' or 'Refused'). The values for 'Don't Know' were specifically filtered out of the listed categories and included here.
  % \end{tablenotes}
\end{table}
An analysis of Table \ref{tab:theft_summary_avg} verifies the trends discussed in \ref{app:longitudinal}. Table \ref{tab:theft_summary_avg} displays that on average, the misuse of existing accounts is the most common form of identity theft. Existing credit card misuse affected the largest share of victims (48.23\%), followed by the fraudulent use of existing bank accounts (41.77\%). These two categories alone demonstrate that the overwhelming majority of incidents involve the compromise of accounts that are already held by the victim. The creation of new fraudulent accounts and the misuse of personal information for other purposes, such as filing a fraudulent tax return, are notably less frequent, impacting 7.82\% and 5.39\% of victims, respectively. Furthermore, a significant portion of the victim population, 13.76\% on average, experienced more than one type of identity theft.

The most common methods by which victims discovered the fraud were being notified by a financial institution/bank (37.69\%), or personally noticing a problem with their account, such as not being able to make purchases anymore (33.65\%), which is a reflection of the prevalence and efficacy of fraud detection algorithms used by financial institutions. On the other hand, the proactive measure of checking credit reports was far less common in discovering misuse, with only 0.81\% of victims reporting that they discovered their incident in this way. Finally, Table \ref{tab:theft_summary_avg} shows that in terms of how the victims' personal information was obtained, no single method is dominant. The most frequently cited pathway was information being stolen during a transaction, either online or in-person, accounting for 12.12\% of incidents. Interestingly, while digital threats are a major public concern, direct hacking or computer theft was reported in only 1.38\% of cases, and data breaches from companies or employers accounted for just 3.73\% on average. It is also noteworthy that traditional, non-digital methods remain relevant. For example, 5.10\% of incidents resulted from a lost or stolen physical item, such as a wallet or mail. However, all these values only reflect the victims who knew how their information was obtained. The vast majority of victims did not know, therefore these statistics may not be an accurate reflection of how the information of most victims was obtained.
%\com{This is purely optional, but some of the stats reporting could be presented in enumeration/bullet form, may make it a bit easier to read..}





\subsection{Summary of Data Demographics}

Recall that Table \ref{tab:victimization_year} showed that the number of identity theft victims has been generally increasing with each survey wave. While identity theft affects all segments of the population, comparison of the data sample's victimization with U.S. Census data reveals some disparities in victimization risk, detailed in Table \ref{tab:demographics_percent} in Appendix \ref{app:demographics_census}. Specifically, victims aged 30-64 are overrepresented among victims. Additionally, white individuals are overrepresented among victims, while  individuals of Black, Asian, American Indian or Alaskan Native descent, or of two or more races, are all underrepresented among victims. On the impact of education level, individuals with a Bachelor's degree or a Graduate/Professional degree are overrepresented among victims. Similar to education, with respect to income, households earning \$75,000 and over are the most victimized group, being slightly overrepresented compared to their share of the population. Households in the lower income brackets (\$25,000 or less) are underrepresented among victims. 

To further investigate why certain groups are disproportionately represented, we conducted a series of deeper analyses, including weighted proportions tests and logistic regressions, extending the work of \cite{Reynolds2021Differential, DeLiema2021} to all six waves of the ITS survey. Our results are detailed in \ref{app:demographics_census} and in the extended version \cite{Anonymous2026SocialCostExtended}, and are summarized as follows. First, we found that victims aged 30-64 are ``digitally active but not digitally native,'' meaning that they are significantly more vulnerable to digital theft than the 16-29 ``digital native'' group. Next, we found disparities in the specific types of fraud experienced across different demographics. Our analysis in the extended version \cite{Anonymous2026SocialCostExtended} reveals that the over-representation of White and highly educated victims is primarily driven by existing credit card misuse, aligning with the results of \cite{Nevin2025, Copes2010}. It was also found that Black individuals and lower-SES groups are disproportionately victimized by more invasive forms of identity theft, such as existing bank account and new account fraud. For example, we found that Black victims report existing bank account misuse at a rate higher than the White victim population, a finding that aligns with prior research on target suitability and financial accessibility \cite{Copes2010}. Finally, the results of our regression analysis in the extended version \cite{Anonymous2026SocialCostExtended} verified the trends described in \cite{Reynolds2021Differential}, where higher levels of age, education, and income increase the odds of identity theft victimization. Thus, the trends observed in previous studies using a fraction of ITS surveys continue to be observed when all six waves are pooled.
%\com{Do any of these observations appear more or less  pronounced/subdued than previously observed, and are there new observations..?}
%\resp{I'm not so sure how to compare them, since their analyses were slightly different than ours, so I can't compare the odds ratios.}

%%%%%%%%%%%%%%%%%%%%%
